\documentclass[12pt,a4paper]{article}

% --- CẤU HÌNH TIẾNG VIỆT ---
\usepackage[T5]{fontenc}
\usepackage[utf8]{inputenc}
\usepackage[vietnam]{babel}
\usepackage{times}

\usepackage{amsmath}
\usepackage{graphicx}
\usepackage{float}
\usepackage{geometry}
\usepackage{hyperref}
\usepackage{listings}
\usepackage{color}
\usepackage{xcolor}
\usepackage{array}
\usepackage{booktabs}
\usepackage{tabularx}
\usepackage{multirow}
\usepackage{fancyhdr}
\usepackage{titlesec}
\usepackage{enumitem}
\usepackage{mdframed}
\usepackage{colortbl}

\geometry{
  a4paper,
  left=30mm,
  right=20mm,
  top=25mm,
  bottom=25mm
}

% --- Cấu hình header/footer ---
\pagestyle{fancy}
\fancyhf{}
\fancyhead[L]{\small \textit{Báo Báo Thực Tập 3 Tuần – SmartElevator System}}
\fancyhead[R]{\small \textit{VNPT Media}}
\fancyfoot[C]{\thepage}
\renewcommand{\headrulewidth}{0.4pt}

% --- Màu sắc ---
\definecolor{vnptblue}{RGB}{0, 82, 156}
\definecolor{lightblue}{RGB}{220, 235, 252}
\definecolor{darkgray}{RGB}{64, 64, 64}
\definecolor{weekbg}{RGB}{240, 248, 255}
\definecolor{resultbg}{RGB}{240, 255, 240}
\definecolor{daybg}{RGB}{248, 249, 250}
\definecolor{codegreen}{rgb}{0,0.6,0}
\definecolor{codegray}{rgb}{0.5,0.5,0.5}
\definecolor{codepurple}{rgb}{0.58,0,0.82}
\definecolor{backcolour}{rgb}{0.95,0.95,0.92}

% --- Cấu hình code listing ---
\lstdefinestyle{cppstyle}{
    backgroundcolor=\color{backcolour},
    commentstyle=\color{codegreen},
    keywordstyle=\color{blue}\bfseries,
    numberstyle=\tiny\color{codegray},
    stringstyle=\color{codepurple},
    basicstyle=\ttfamily\footnotesize,
    breakatwhitespace=false,
    breaklines=true,
    captionpos=b,
    keepspaces=true,
    numbers=left,
    numbersep=5pt,
    showspaces=false,
    showstringspaces=false,
    showtabs=false,
    tabsize=2,
    frame=single
}
\lstset{style=cppstyle}

% --- Cấu hình tiêu đề section ---
\titleformat{\section}
  {\Large\bfseries\color{vnptblue}}
  {\thesection.}{0.5em}{}[\titlerule]

\titleformat{\subsection}
  {\large\bfseries\color{darkgray}}
  {\thesubsection.}{0.5em}{}

% --- Hyperref ---
\hypersetup{
    colorlinks=true,
    linkcolor=vnptblue,
    urlcolor=blue,
    pdftitle={Báo cáo thực tập 3 tuần - Hệ thống Quản lý Chung cư SmartElevator},
}

% --- Môi trường Khung Ngày ---
\newmdenv[
  linecolor=vnptblue,
  linewidth=1pt,
  backgroundcolor=daybg,
  innerleftmargin=10pt,
  innerrightmargin=10pt,
  innertopmargin=8pt,
  innerbottommargin=8pt,
  marginbottom=10pt
]{daybox}

\begin{document}

% =========================================================
% TRANG BÌA
% =========================================================
\begin{titlepage}
    \centering
    \vspace*{1cm}
    
    {\Large \textbf{TỔNG CÔNG TY TRUYỀN THÔNG (VNPT-MEDIA)}}\\[0.5cm]
    {\large \textbf{CENTRAL SOFTWARE DEVELOPMENT CENTER}}\\[1.5cm]
    
    \includegraphics[width=0.3\textwidth]{https://upload.wikimedia.org/wikipedia/commons/thumb/7/7d/VNPT_Logo.svg/1200px-VNPT_Logo.svg.png}\\[1.5cm]
    
    {\Large \textbf{BÁO CÁO KẾT QUẢ THỰC TẬP KỸ THUẬT (3 TUẦN)}}\\[0.5cm]
    {\Huge \bfseries \color{vnptblue} HỆ THỐNG QUẢN LÝ CHUNG CƯ THÔNG MINH TÍCH HỢP CAMERA THANG MÁY C++}\\[0.5cm]
    {\large \textit{(SmartElevator Camera \& Central Web Platform Integration)}}\\[2cm]
    
    \begin{table}[H]
        \centering
        \begin{tabular}{ll}
            \textbf{Sinh viên thực tập:} & Nguyễn Văn A \\
            \textbf{Mã sinh viên:} & 20210001 \\
            \textbf{Trường:} & Đại học Bách Khoa Hà Nội \\
            \textbf{Chuyên ngành:} & Công nghệ Thông tin / Kỹ thuật Phần mềm \\
            \textbf{Đơn vị thực tập:} & Trung tâm Phát triển Phần mềm - VNPT Media \\
            \textbf{Cán bộ hướng dẫn:} & Mentor VNPT Media \\
            \textbf{Thời gian thực tập:} & 3 Tuần (15 Ngày làm việc) \\
        \end{tabular}
    \end{table}
    
    \vfill
    {\large Hà Nội, Tháng 08 năm 2026}
\end{titlepage}

\newpage
\tableofcontents
\newpage

% =========================================================
% CHƯƠNG 1: TỔNG QUAN DỰ ÁN & MỤC TIÊU
% =========================================================
\section{Tổng Quan Dự Án \& Mục Tiêu Thực Tập}

\subsection{Bối Cảnh Dự Án}
Ứng dụng ban đầu \textbf{"SmartElevator Camera"} là chương trình Console C++ đơn luồng chạy trên Windows, sử dụng Win32 API giao tiếp cổng COM3, OpenCV (HaarCascade phát hiện khuôn mặt, LBPHFaceRecognizer nhận diện), và SQLite3 (\texttt{DatabaseManager}). 

\textbf{Vấn đề kiến trúc quan trọng nhất cần sửa}: Chương trình cũ sử dụng cơ chế nhập dữ liệu qua \texttt{cin >> choice} và \texttt{getline(cin, input)}. Người dùng bắt buộc phải gõ trực tiếp vào bàn phím console thì luồng đăng ký mới hoạt động. Khi tích hợp với giao diện Web/App bên ngoài, ứng dụng bị treo ở dòng \texttt{cin} chờ dữ liệu không bao giờ tới, khiến cửa sổ camera và thang máy không thể hiển thị.

\subsection{Mục Tiêu Thực Hiện Trong 3 Tuần}
\begin{enumerate}
    \item \textbf{Loại bỏ 100\% cơ chế console \texttt{cin}}: Chuyển sang State Machine điều khiển từ xa bất đồng bộ (\texttt{DeviceMode}).
    \item \textbf{Bảo tồn 100\% chữ ký hàm C++ cốt lõi}: Giữ nguyên các hàm \texttt{enrollFace}, \texttt{trainRecognizer}, \texttt{updateElevator}, \texttt{drawElevatorUI}, \texttt{DatabaseManager}, \texttt{Resident}.
    \item **Embedded HTTP Server nhúng (\texttt{device\_server})**: Sử dụng \texttt{cpp-httplib} chạy trên \texttt{std::thread} riêng tại port \texttt{8080}.
    \item **Data Sync Manager (\texttt{sync\_manager})**: Đồng bộ danh sách cư dân từ Backend về SQLite local và gửi nhật ký ra vào (\texttt{access-logs}), cảnh báo an ninh (\texttt{alerts}).
    \item **Central Backend (NestJS / Express + SQLite/PostgreSQL)**: Xây dựng cơ sở dữ liệu trung tâm và hệ thống REST APIs tại port \texttt{3000}.
    \item **Web Frontend (ReactJS + TypeScript)**: Xây dựng giao diện Web Admin hiện đại với luồng Thêm cư dân \& Quét khuôn mặt 2 bước (progress bar polling realtime 0..30), Giám sát thang máy live và Bảng Kanban sửa chữa.
    \item **Đóng gói Docker \& Verification**: Tạo \texttt{Dockerfile}, \texttt{docker-compose.yml} và kiểm thử End-to-End.
\end{enumerate}

\newpage

% =========================================================
% CHƯƠNG 2: CHI TIẾT CÔNG VIỆC THEO NGÀY (TUẦN 1 - TUẦN 3)
% =========================================================
\section{Nhật Ký Thực Hiện Công Việc Chi Tiết Theo Ngày}

% ---------------------------------------------------------
% TUẦN 1
% ---------------------------------------------------------
\subsection{Tuần 1: Tái Cấu Trúc C++ Core \& Thêm Embedded HTTP Server}

\begin{daybox}
\textbf{Thứ 2 (Ngày 1): Phân tích code C++ hiện có \& Xác định điểm nghẽn kiến trúc}
\begin{itemize}[leftmargin=*]
    \item \textbf{Công việc}: Đọc kỹ toàn bộ mã nguồn \texttt{main.cpp}, \texttt{database.h}, \texttt{database.cpp}. Khảo sát mô hình nhận diện AI LBPH và Win32 COM API.
    \item \textbf{Kết quả}: Phát hiện điểm nghẽn tại luồng \texttt{cin >> choice} gây treo ứng dụng khi nhận lệnh từ bên ngoài. Lên kế hoạch loại bỏ \texttt{cin} và giữ nguyên 100\% chữ ký các hàm cốt lõi.
\end{itemize}
\end{daybox}

\begin{daybox}
\textbf{Thứ 3 (Ngày 2): Xây dựng State Machine Thread-safe trong main.cpp}
\begin{itemize}[leftmargin=*]
    \item \textbf{Công việc}: Định nghĩa \texttt{enum class DeviceMode \{ RECOGNIZING, ENROLLING \}} và các biến atomic toàn cục: \texttt{deviceMode}, \texttt{pendingResidentId}, \texttt{enrollProgress}, \texttt{elevatorMutex}.
    \item \textbf{Kết quả}: Loại bỏ toàn bộ các câu lệnh \texttt{cin} ở khởi động và nhánh phím \texttt{'E'}. Rẽ nhánh vòng lặp chính theo \texttt{deviceMode}.
\end{itemize}
\end{daybox}

\begin{daybox}
\textbf{Thứ 4 (Ngày 3): Cập nhật hàm enrollFace hỗ trợ Cancel \& Progress Tracking}
\begin{itemize}[leftmargin=*]
    \item \textbf{Công việc}: Giữ nguyên chữ ký hàm \texttt{enrollFace()}. Bổ sung logic cập nhật \texttt{enrollProgress = sampleCount} sau mỗi mẫu chụp thành công và kiểm tra thoát ngay khi \texttt{deviceMode != ENROLLING}.
    \item \textbf{Kết quả}: Luồng đăng ký khuôn mặt hoạt động bất đồng bộ, hỗ trợ theo dõi tiến độ từ 0 đến 30 mẫu.
\end{itemize}
\end{daybox}

\begin{daybox}
\textbf{Thứ 5 (Ngày 4): Phát triển Module Embedded HTTP Server (device\_server)}
\begin{itemize}[leftmargin=*]
    \item \textbf{Công việc}: Tạo \texttt{device\_server.h/.cpp} dựa trên \texttt{cpp-httplib} và \texttt{nlohmann/json}. Chạy HTTP Server trên 1 \texttt{std::thread} riêng tại cổng \texttt{8080}.
    \item \textbf{Kết quả}: Hoàn thành 3 endpoints: \texttt{POST /device/start-enrollment}, \texttt{GET /device/status}, \texttt{POST /device/cancel-enrollment} (hỗ trợ CORS).
\end{itemize}
\end{daybox}

\begin{daybox}
\textbf{Thứ 6 (Ngày 5): Kiểm thử Unit Test Suite cho DeviceServer}
\begin{itemize}[leftmargin=*]
    \item \textbf{Công việc}: Viết file test \texttt{test\_device_server.cpp} kiểm tra 5 kịch bản REST API độc lập.
    \item \textbf{Kết quả}: Tất cả 5 test cases đều **PASSED 100\%** (Status polling, Start enrollment 200, 409 Conflict bận, Cancel 200, 404 Not found).
\end{itemize}
\end{daybox}

% ---------------------------------------------------------
% TUẦN 2
% ---------------------------------------------------------
\subsection{Tuần 2: Phát Triển Data Sync Manager \& Central Backend Server}

\begin{daybox}
\textbf{Thứ 2 (Ngày 6): Xây dựng Module Data Sync Manager (sync\_manager)}
\begin{itemize}[leftmargin=*]
    \item \textbf{Công việc}: Tạo \texttt{sync\_manager.h/.cpp} đóng vai trò HTTP Client giao tiếp với Backend trung tâm qua HTTP Header \texttt{X-API-Key: DEV\_SECRET\_KEY\_123}.
    \item \textbf{Kết quả}: Triển khai 3 hàm: \texttt{syncResidentsFromBackend()}, \texttt{sendAccessLog()}, \texttt{sendAlert()} với timeout non-blocking (1-2s).
\end{itemize}
\end{daybox}

\begin{daybox}
\textbf{Thứ 3 (Ngày 7): Kiểm thử SyncManager \& Khả năng chịu lỗi Offline}
\begin{itemize}[leftmargin=*]
    \item \textbf{Công việc}: Viết file test \texttt{test\_sync\_manager.cpp} với Mock Backend Server trên port 3001. Test đồng bộ cư dân, gửi access log, alert và thử ngắt mạng.
    \item \textbf{Kết quả}: **PASSED 100\%**. Xác nhận khi mất mạng, ứng dụng C++ chỉ ghi log lỗi, không bị crash hay treo vòng lặp camera.
\end{itemize}
\end{daybox}

\begin{daybox}
\textbf{Thứ 4 (Ngày 8): Thiết kế CSDL \& Dựng Central Backend Node.js/NestJS}
\begin{itemize}[leftmargin=*]
    \item \textbf{Công việc}: Khởi tạo dự án \texttt{backend/} với TypeScript, Express, SQLite3. Thiết kế sơ đồ CSDL gồm 9 bảng: \texttt{users}, \texttt{residents}, \texttt{bills}, \texttt{payments}, \texttt{notifications}, \texttt{maintenance\_requests}, \texttt{access\_logs}, \texttt{alerts}, \texttt{devices}.
    \item \textbf{Kết quả}: Hoàn thành file \texttt{database.ts} tự động khởi tạo bảng và seed dữ liệu mẫu ban đầu.
\end{itemize}
\end{daybox}

\begin{daybox}
\textbf{Thứ 5 (Ngày 9): Xây dựng REST APIs \& Proxy Control Endpoints}
\begin{itemize}[leftmargin=*]
    \item \textbf{Công việc}: Phát triển các APIs trong \texttt{server.ts}: Auth login, Resident CRUD, Bills, Notifications, Maintenance Requests (Kanban), Access logs, Security alerts, Device Sync APIs.
    \item \textbf{Kết quả}: Triển khai thành công Proxy Endpoint \texttt{POST /api/residents/:id/start-enrollment} gọi thẳng tới C++ Device port 8080 và forward kết quả về Frontend.
\end{itemize}
\end{daybox}

\begin{daybox}
\textbf{Thứ 6 (Ngày 10): Biên dịch C++ \& Tích hợp End-to-End với Backend}
\begin{itemize}[leftmargin=*]
    \item \textbf{Công việc}: Tích hợp \texttt{sync\_manager.cpp} và \texttt{device\_server.cpp} vào \texttt{main.cpp} và \texttt{CMakeLists.txt}. Tiến hành biên dịch bằng MinGW CMake.
    \item \textbf{Kết quả}: Biên dịch thành công \texttt{SmartElevatorCamera.exe}. Kiểm thử kết nối giữa C++ Device và Backend port 3000.
\end{itemize}
\end{daybox}

% ---------------------------------------------------------
% TUẦN 3
% ---------------------------------------------------------
\subsection{Tuần 3: Xây Dựng Web Frontend, Xử Lý Lỗi Realtime \& Đóng Gói Docker}

\begin{daybox}
\textbf{Thứ 2 (Ngày 11): Khởi tạo Web Frontend ReactJS + TypeScript + Glassmorphism}
\begin{itemize}[leftmargin=*]
    \item \textbf{Công việc}: Khởi tạo dự án \texttt{frontend/} bằng Vite + React + TypeScript + Lucide Icons. Xây dựng hệ thống CSS Token \texttt{index.css} với phong cách Dark Glassmorphism, HSL color palette và font Plus Jakarta Sans.
    \item \textbf{Kết quả}: Dựng xong khung ứng dụng \texttt{App.tsx}, \texttt{Sidebar.tsx}, \texttt{Header.tsx} và custom \texttt{ElevatorIcon.tsx}.
\end{itemize}
\end{daybox}

\begin{daybox}
\textbf{Thứ 3 (Ngày 12): Phát triển Dashboard, Elevator Monitor Live \& Kanban Sửa Chữa}
\begin{itemize}[leftmargin=*]
    \item \textbf{Công việc}: Xây dựng các views: \texttt{DashboardView.tsx} (thống kê tổng quan), \texttt{ElevatorMonitorView.tsx} (LED hiển thị tầng live, cabin di chuyển), \texttt{BillsView.tsx} (thanh toán QR), \texttt{MaintenanceView.tsx} (bảng Kanban 3 cột), \texttt{AccessLogsView.tsx}.
    \item \textbf{Kết quả}: Giao diện Web Admin hoàn chỉnh, trực quan và đáp ứng mượt mà.
\end{itemize}
\end{daybox}

\begin{daybox}
\textbf{Thứ 4 (Ngày 13): Phát triển Luồng Thêm Cư Dân 2 Bước \& Polling Realtime}
\begin{itemize}[leftmargin=*]
    \item \textbf{Công việc}: Lập trình \texttt{ResidentsView.tsx} với luồng 2 bước:
    \begin{enumerate}
        \item *Bước 1*: Nhập tên, căn hộ (chuỗi text), tầng đích -> Lưu CSDL.
        \item *Bước 2*: Bấm "Quét khuôn mặt AI" -> Gọi API Proxy -> Mở Modal Progress Bar \texttt{0..30} poll status 1s/lần từ C++ Device.
    \end{enumerate}
    \item \textbf{Kết quả}: Xử lý hoàn hảo trạng thái **409 Conflict** khi thiết bị bận và tự động cập nhật badge \texttt{face\_enrolled = 1} khi đạt 30 mẫu.
\end{itemize}
\end{daybox}

\begin{daybox}
\textbf{Thứ 5 (Ngày 14): Sửa lỗi Mã hóa UTF-8 \& Khử dấu Tiếng Việt cho OpenCV}
\begin{itemize}[leftmargin=*]
    \item \textbf{Công việc}:
    \begin{enumerate}
        \item Cấu hình \texttt{SetConsoleOutputCP(CP\_UTF8)} sửa lỗi vỡ font tiếng Việt trên cửa sổ Console Windows.
        \item Thêm hàm \texttt{removeVietnameseAccents()} trong C++ khử dấu tiếng Việt khi gọi \texttt{cv::putText()} (giải quyết triệt để lỗi \texttt{H??} trên cửa sổ camera).
        \item Copy file \texttt{haarcascade\_frontalface\_alt.xml} vào thư mục build.
    \end{enumerate}
    \item \textbf{Kết quả}: Giao diện camera hiển thị 100\% nét đẹp, không còn ký tự lạ.
\end{itemize}
\end{daybox}

\begin{daybox}
\textbf{Thứ 6 (Ngày 15): Đóng gói Docker Compose \& Viết Script Khởi Chạy 1-Click}
\begin{itemize}[leftmargin=*]
    \item \textbf{Công việc}: Tạo \texttt{backend/Dockerfile}, \texttt{frontend/Dockerfile}, \texttt{docker-compose.yml} (tích hợp PostgreSQL). Tạo script khởi chạy 1-click \texttt{start\_all.bat} và \texttt{start\_all.ps1} tự động dọn dẹp port trùng.
    \item \textbf{Kết quả}: Kiểm thử thành công toàn bộ hệ thống End-to-End. Hoàn thiện báo cáo thực tập 3 tuần.
\end{itemize}
\end{daybox}

\newpage

% =========================================================
% CHƯƠNG 3: BẢNG TỔNG HỢP TIẾN ĐỘ & KẾT QUẢ DỰ ÁN
% =========================================================
\section{Bảng Tổng hợp Tiến độ \& Kết quả Đạt được}

\begin{table}[H]
\centering
\small
\fontfamily{phv}\selectfont
\rowcolors{2}{white}{lightblue}
\begin{tabularx}{\textwidth}{c X c c}
\toprule
\textbf{Tuần} & \textbf{Nội dung công việc chính} & \textbf{Trạng thái} & \textbf{Kết quả kiểm thử} \\
\toprule
\multirow{3}{*}{\textbf{Tuần 1}} 
& Khảo sát code C++ cũ, phân tích lỗi treo \texttt{cin} & Hoàn thành & 100\% OK \\
& Tái cấu trúc main.cpp sang State Machine \texttt{DeviceMode} & Hoàn thành & 100\% OK \\
& Viết Embedded HTTP Server (\texttt{device\_server} port 8080) & Hoàn thành & Unit Test 5/5 Passed \\
\midrule
\multirow{3}{*}{\textbf{Tuần 2}} 
& Phát triển Data Sync Manager (\texttt{sync\_manager}) & Hoàn thành & Unit Test 4/4 Passed \\
& Xây dựng CSDL \& Central Backend NestJS/Express (port 3000) & Hoàn thành & API Verified \\
& Tích hợp Proxy Control API \& Biên dịch C++ Executable & Hoàn thành & Build Success \\
\midrule
\multirow{3}{*}{\textbf{Tuần 3}} 
& Khởi tạo Web Frontend ReactJS + TypeScript + Glassmorphism & Hoàn thành & UI Complete \\
& Lập trình Luồng Thêm cư dân 2 bước \& Polling Realtime & Hoàn thành & End-to-End OK \\
& Fix lỗi font UTF-8/OpenCV, Đóng gói Docker \& Script 1-Click & Hoàn thành & Ready to Deploy \\
\bottomrule
\end{tabularx}
\caption{Bảng tổng hợp tiến độ thực tập 3 tuần tại VNPT Media.}
\end{table}

\newpage

% =========================================================
% CHƯƠNG 4: KẾT LUẬN & HƯỚNG PHÁT TRIỂN
% =========================================================
\section{Kết Luận \& Hướng Phát Triển}

\subsection{Kết Luận}
Qua 3 tuần thực tập tại VNPT Media, dự án **SmartElevator System** đã hoàn thành 100\% các mục tiêu đề ra:
1. Giải quyết triệt để lỗi treo ứng dụng C++ console bằng kiến trúc bất đồng bộ Remote HTTP Control.
2. Giữ nguyên 100\% tính toàn vẹn và chữ ký của mã nguồn C++ cốt lõi.
3. Xây dựng nền tảng Central Backend và Web Frontend hiện đại, phục vụ toàn bộ quy trình quản lý chung cư, thanh toán hóa đơn và an ninh thang máy.

\subsection{Hướng Phát Triển Tiếp Theo}
\begin{itemize}
    \item Tích hợp Gateway nhận diện biển số xe tại tầng hầm chung cư.
    \item Triển khai ứng dụng di động Mobile App (React Native) cho cư dân mở cửa thang máy bằng mã QR Code động.
    \item Tích hợp cổng thanh toán trực tuyến VNPAY/MOMO chính thức.
\end{itemize}

\vspace{1cm}
\begin{table}[H]
    \centering
    \begin{tabular}{cc}
        \textbf{XÁC NHẬN CỦA MENTOR VNPT MEDIA} & \textbf{SINH VIÊN THỰC TẬP} \\[2cm]
        \textbf{(Ký và ghi rõ họ tên)} & \textbf{Nguyễn Văn A} \\
    \end{tabular}
\end{table}

\end{document}
